\documentclass[10pt,a4paper]{amsart}
\usepackage[margin=19mm]{geometry}
\usepackage{amsmath,amssymb,mathtools}
\usepackage[scr=rsfs,cal=euler]{mathalfa}
\usepackage{xfrac}
\usepackage[unicode,hidelinks,bookmarks=false]{hyperref}
\newcommand{\ie}{i.e.\ }
\newcommand{\cf}{cf.\ }
\newcommand{\resp}{respectively\ }
\newcommand{\Subsection}{\subsection}
\providecommand{\nobreakdash}{\nobreak\hbox{-}}
\setlength{\emergencystretch}{2em}
\multlinegap0pt
\usepackage{bm}
\usepackage{bbm}
\usepackage{dsfont}
\usepackage{xspace}
\usepackage{cancel}
\usepackage{mathtools}
\usepackage{amsthm}
\makeatletter
\@ifundefined{theorem}{\newtheorem{theorem}{Theorem}}{}
\@ifundefined{lemma}{\newtheorem{lemma}[theorem]{Lemma}}{}
\@ifundefined{proposition}{\newtheorem{proposition}[theorem]{Proposition}}{}
\makeatother
\title[Real Reliability Roots of Simple Graphs are Dense]{Real Reliability Roots of Simple Graphs are Dense}
\author{Mohamed Omar}
\date{Source: arXiv:2604.03530v1 (CC BY 4.0)}
\begin{document}
\maketitle

\noindent\emph{Source and licence.} Real Reliability Roots of Simple Graphs are Dense. Version arXiv:2604.03530v1 (CC BY 4.0), distributed under the Creative Commons Attribution 4.0 International licence, as recorded in the arXiv licence metadata for this exact version and shown on the abstract page https://arxiv.org/abs/2604.03530v1. Full rights notice: licenses/arxiv-2604.03530-CC-BY-4.0.txt.\par\smallskip

\noindent\textit{Editorial scope.} Counted unit: the Proposition whose source label is \texttt{prop:yn} in the article (source0.tex lines 246--248, source label \texttt{prop:yn}), delivered together with the article's own complete original proof of it (source0.tex lines 275--374, one \texttt{proof} environment of 100 lines). No mathematics is written, repaired or re-derived: statement and proof are the article's own text line for line. The proof is detached from the statement in the source: the article states the result at lines 246--248 and prints its proof at lines 275--374, and the proof environment's own header names the statement, so the association is the article's own. Required context retained uncounted: prop:recurrences (lines 93--105); eq:complete-rec (lines 95--97); eq:split-rec (lines 98--100); lem:kernel (lines 204--213). Every definition, display and label of those lines is retained unchanged. Omitted from this package and not redistributed: 1--92, 101--203, 214--245, 249--274, 375--494. Nothing inside a retained excerpt is omitted or altered: every retained body is delivered exactly as the article's lines are written, so no omission receipt of any kind accompanies this package. Citations: the retained excerpts carry no citation command at all, so no bibliography entry of the article is transcribed and no reference is redistributed. Reference adaptation: none was needed and none was made; every reference inside the delivered unit resolves inside it, so no unresolved reference or citation is produced. Printed slips kept literal: no printed word, symbol, hypothesis or number of the counted statement or of its proof was altered. Layout and class: the article's own class is \texttt{amsart} with packages including geometry, amsmath, amssymb, amsfonts, mathtools, hyperref, tikz, bbm, xcolor; the wrapper is the lane's pinned \texttt{amsart} editorial wrapper at 10pt on A4, which restates the theorem-like environments the retained text needs and adds the article's own macro definitions that the retained text needs (none), with extra wrapper packages (bm, bbm, dsfont, xspace, cancel, mathtools). The delivered numbering is the unit's own. Assets: no retained span contains a figure environment, a graphics-inclusion command, a PDF-page inclusion, a table or a TikZ picture; no image file is copied into this package, the package declares no asset dependency and no image file is redistributed with it. Licence: version arXiv:2604.03530v1 of this article is distributed under the Creative Commons Attribution 4.0 International licence, as recorded in the arXiv licence metadata for this exact version and shown on the abstract page https://arxiv.org/abs/2604.03530v1. A full rights notice is delivered at licenses/arxiv-2604.03530-CC-BY-4.0.txt. Characters: every retained excerpt is pure ASCII, so the delivered engine, XeLaTeX with its default Latin Modern text font, reports no missing character.
\begin{proposition}\label{prop:recurrences}
The following recurrences hold for $n \geq 2$:
\begin{equation}\label{eq:complete-rec}
C_n(q)=1-\sum_{a=1}^{n-1} \binom{n-1}{a-1}C_a(q)q^{a(n-a)},   
\end{equation}
\begin{equation}\label{eq:split-rec}
S_n(q)=\sum_{a=1}^{n-1} \binom{n-2}{a-1} C_a(q)C_{n-a}(q)q^{a(n-a)-1},
\end{equation}
\begin{equation}\label{eq:rminus-rec}
R_n(q)=1 - \sum_{a=1}^{n-1} \binom{n-2}{a-1} C_a(q) q^{a(n-a)-1} - \sum_{a=3}^{n-1} \binom{n-2}{a-2} R_a(q) q^{a(n-a)}.
\end{equation}
Furthermore $C_1(q)=1$.
\end{proposition}

\begin{lemma}\label{lem:kernel}
Let $K \subset (-1,0)$ be compact. Fix an integer $s$ and let $b_{N,a}(q)$ be real-valued functions indexed by natural numbers $a,N$ with $N \geq 2$ and $1 \leq a \leq N-1$. Suppose
\[
\sup \ \{ |b_{N,a}(q)| \colon N \geq 2, \ 1 \leq a \leq N-1, \ q \in K \} <\infty.
\]
Then
\[
\sup_{q \in K} \left| \sum_{a=1}^{N-1} \binom{N}{a}b_{N,a}(q) \cdot q^{a(N-a)+s} \right| \to 0 \qquad \text{ as } N \to \infty.
\]
\end{lemma}

\begin{proposition}\label{prop:yn}
Let $K \subset (-1,0)$ be compact. Then there exists an integer $N$ such that $\hat{y}_N(q)<-1$ for every $q \in K$.
\end{proposition}

\begin{proof}[Proof of Proposition~\ref{prop:yn}]
The proof has three parts. First, we use Proposition~\ref{prop:recurrences} and Lemma~\ref{lem:kernel} to prove the families of functions $\{C_n\}_{n \geq 1}$ and $\{R_n\}_{n \geq 1}$ are uniformly bounded on $K$. Then, we use this to show that uniformly on $K$, 
\[
C_n(q) \to 1, \qquad R_n(q) \to 1, \qquad q^{2-n}S_n(q) \to 2.
\]
It will then follow that $2q^{n-2}\hat{y}_n(q) \to 1$ uniformly on $K$ and this will lead to a sufficiently large $N$ for which $\hat{y}_N(q)<-1$ for all $q \in K$.

We start by showing $\{C_n\}_{n \geq 1},\{R_n\}_{n \geq 1}$ are uniformly bounded on $K$. For $n \geq 2$, let
\[
X_n:=\sup_{q \in K} |C_n(q)|, \qquad Y_n:=\sup_{q \in K} |R_n(q)|, \qquad M_n:=\max\{X_k,Y_k \colon 1 \leq k \leq n\}.
\]
From \eqref{eq:complete-rec} of Proposition~\ref{prop:recurrences}, and the fact that $\binom{n-1}{a-1}=\binom{n}{a}\frac{a}{n}$, we have that for any $q \in K$,
\[
|C_n(q)| \leq 1 + M_{n-1} \sum_{a=1}^{n-1} \binom{n}{a} \frac{a}{n} |q|^{a(n-a)}=1+M_{n-1}\sum_{a=1}^{n-1} \binom{n}{a} \frac{a}{n} (-1)^{a(n-a)} q^{a(n-a)},
\]
the latter equality because $q<0$. Taking the supremum over $q \in K$ we have $X_n \leq 1+M_{n-1}\alpha_n$ where 
\[
\alpha_n=\sup_{q \in K} \sum_{a=1}^{n-1} \binom{n}{a} \frac{a}{n} (-1)^{a(n-a)} q^{a(n-a)}.
\] 
In $\alpha_n$ we can replace the sum by its absolute value because the sum is positive since $q<0$. Therefore, we can apply  Lemma~\ref{lem:kernel} with $b_{n,a}(q)=\frac{a}{n}(-1)^{a(n-a)}$, $N=n$ and $s=0$. We see that $|b_{n,a}(q)|$ is uniformly bounded by $1$ for all $n,a$. Therefore $\alpha_n \to 0$ as $n \to \infty$.

Applying a similar process, with the observation that 
\[
\binom{n-2}{a-1}=\frac{a(n-a)}{n(n-1)}\binom{n}{a} \qquad{ \text{and} } \qquad \binom{n-2}{a-2}=\frac{a(a-1)}{n(n-1)}\binom{n}{a}
\]
we have $Y_n \leq 1+M_{n-1}\beta_n$ where
\[
\beta_n := \sup_{q \in K} \sum_{a=1}^{n-1} \binom{n}{a}\frac{a(n-a)}{n(n-1)}(-1)^{a(n-a)-1} q^{a(n-a)-1} + \sup_{q \in K} \sum_{a=3}^{n-1} \binom{n}{a}\frac{a(a-1)}{n(n-1)}(-1)^{a(n-a)} q^{a(n-a)}. 
\]
Apply Lemma~\ref{lem:kernel} to the first of the two sums with $N=n,s=-1,b_{n,a}(q)=\frac{a(n-a)}{n(n-1)}(-1)^{a(n-a)}$, and the same lemma to the second sum with $N=n,s=0,b_{n,a}(q)=\frac{a(a-1)}{n(n-1)}(-1)^{a(n-a)}$ for $3 \leq a \leq n-1$ and $b_{n,a}(q)=0$ for $1 \leq a \leq 2$, this yields $\beta_n \to 0$ as $n \to \infty$.

For any $n$ let $\gamma_n=\max\{\alpha_n,\beta_n\}.$ Then we have $\gamma_n\to 0$. Furthermore, for $n\ge 2$, we have $X_n\le 1+M_{n-1}\gamma_n$ and $Y_n\le 1+M_{n-1}\gamma_n$, so $M_n=\max\{M_{n-1},X_n,Y_n\}\le \max\{M_{n-1},\,1+M_{n-1}\gamma_n\}.$ Choose $N \geq 2$ such that \(\gamma_n\le \tfrac12\) for all \(n\ge N\). Then for all \(n\ge N\),
\[
M_n\le \max\left\{M_{n-1},\,1+\frac12 M_{n-1}\right\}.
\]
Set
\[
C=\max\{M_1,\dots,M_{N-1},2\}.
\]
We prove by induction that $M_n\le C$ for all $n$. This is immediate for $n\le N-1$. If $n\ge N$ and $M_{n-1}\le C$, then
\[
X_n\le 1+\frac12 M_{n-1}\le 1+\frac12 C\le C,
\]
the last inequality because $C \geq 2$. Similarly $Y_n\le C$. Therefore $M_n=\max\{M_{n-1},X_n,Y_n\}\le C.$ Thus $M_n\le C$ for all $n$. Since $0\le X_n,Y_n\le M_n$, it follows that the polynomials $\{C_n\}_{n \geq 1},\{R_n\}_{n \geq 1}$ are uniformly bounded on $K$.

Next we show that
\[
C_n(q) \to 1, \qquad R_n(q) \to 1, \qquad q^{2-n}S_n(q) \to 2,
\]
uniformly on $K$ as $n \to \infty$. From Proposition~\ref{prop:recurrences},
\[
C_n(q)-1=-\sum_{a=1}^{n-1}\binom{n}{a}\frac{a}{n}C_a(q)q^{a(n-a)}.
\]
The coefficients $-\frac{a}{n}C_a(q)$ are bounded in absolute value by $C$, so Lemma~\ref{lem:kernel} implies $|C_q(n)-1| \to 0$ uniformly on $K$, and hence $C_q(n) \to 1$ uniformly on $K$. Similarly,
\[
R_n(q)-1 = -\sum_{a=1}^{n-1} \binom{n}{a} \frac{a(n-a)}{n(n-1)}C_a(q)q^{a(n-a)-1} - \sum_{a=3}^{n-1} \binom{n}{a} \frac{a(a-1)}{n(n-1)}R_a(q)q^{a(n-a)}.
\]
The coefficients $-\frac{a(n-a)}{n(n-1)}C_a(q)$ and $-\frac{a(a-1)}{n(n-1)}R_a(q)$ are both bounded above by $C$ in absolute value, so applying Lemma~\ref{lem:kernel} with $s=-1$ to the first sum and $s=0$ to the second we get $R_n(q) \to 1$ uniformly on $K$. Finally, to show $q^{2-n}S_n(q) \to 2$, start by isolating the $a=1$ and $a=n-1$ terms in the sum \eqref{eq:split-rec} in Proposition~\ref{prop:recurrences} and multiply by $q^{2-n}$. Together with the fact $C_1(q)=1$, this yields,
\begin{align*}
q^{2-n}S_n(q) &= 2C_{n-1}(q) + \sum_{a=2}^{n-2} \binom{n-2}{a-1} C_a(q)C_{n-a}(q)q^{a(n-a)-1+2-n} \\
&= 2C_{n-1}(q) + \sum_{a=2}^{n-2} \binom{n-2}{a-1} C_a(q)C_{n-a}(q)q^{(a-1)(n-a-1)}.
\end{align*}
Reindexing the sum with $N=n-2$ and $b=1$, we have 
\[
\left|q^{2-n}S_n(q)-2C_{n-1}(q)\right|=\left|\sum_{b=1}^{N-1} \binom{N}{b} C_{b+1}(q)C_{N+1-b}(q)q^{b(N-b)}\right|.
\]
By Lemma~\ref{lem:kernel}, this vanishes because $|C_{b+1}(q)C_{N+1-b}(q)|$ is uniformly bounded by $C^2$ on $K$. Since $2C_{n-1}(q)$ converges to $1$ uniformly on $K$, it follows $q^{2-n}S_n(q) \to 2$ uniformly on $K$. 

We now find a positive integer $N$ such that $\hat{y}_N(q)<-1$ for all $q \in K$. Since $q^{2-n}S_n(q) \to 2$ uniformly on $K$, we can select $N_0$ so that for all $q \in K$
\[
|q^{2-n}S_n(q)-2|<1 \qquad \forall n \geq N_0.
\]
For such $n$, $q^{2-n}S_n(q)$, and hence $S_n(q)$, are nonzero. So, $\hat{y}_n(q)$ is defined. We then have
\begin{align*}
|2q^{n-2}\hat{y}_n(q)-1| &=\left|2q^{n-2} \cdot \left(\frac{R_n(q)}{S_n(q)}+1\right)-1\right|\\
&= \left| \frac{2R_n(q)}{q^{2-n}S_n(q)} - 1 + 2q^{n-2} \right| \\
&\leq \frac{|2R_n(q)-q^{2-n}S_n(q)|}{|q^{2-n}S_n(q)|} + 2|q|^{n-2} \\
&\leq |2R_n(q)-q^{2-n}S_n(q)|+2|q|^{n-2} \\
&\leq |2R_n(q)-2|+|2-q^{2-n}S_n(q)|+2|q|^{n-2},
\end{align*}
the second-to-last inequality holding because $|q^{2-n}S_n(q)-2|<1$ for $n \geq N_0$ and $q \in K$ implies $q^{2-n}S_n(q)>1$ for the same values of $n$ and $q$. Now $2R_n(q) \to 2$ and $q^{n-2}S_n(q) \to 2$ uniformly on $K$. Moreover since $K$ is a compact subset of $(-1,0)$, there is a universal $\rho \in (0,1)$ so that $|q|\leq \rho$ for $q \in K$ so $2|q|^{n-2} \to 0$ uniformly on $K$. Consequently
\[
2q^{n-2}\hat{y}_n(q) \to 1
\]
uniformly on $K$. So there exists $N_1$ such that for all $q \in K$,
\[
|2q^{n-2}\hat{y}_n(q)-1|<\frac{1}{2} \qquad \forall n \geq N_1.
\]
Moreover, there exists $N_2$ such that for $n \geq N_2$, 
\[
\rho^{n-2}<\frac{1}{4} \qquad \forall n \geq N_2,
\]
because $0<\rho<1$. 

Select an odd integer $N \geq 3$ so that $N \geq \max(N_0,N_1,N_2)$. We claim that if $q \in K$ then $\hat{y}_N(q)<-1$. Indeed select $q \in K$. Since $N \geq N_0,N_1$, we have $2q^{N-2}\hat{y}_N(q)>1/2$. Now since $q<0$ and $N$ is odd, $q^{N-2}<0$. Therefore
\[
\hat{y}_N(q)<\frac{1}{4q^{N-2}} \leq \frac{-1}{4\rho^{N-2}}.
\]
Finally, $N \geq N_2$ implies $\hat{y}_N(q)<-1$. 
\end{proof}

\end{document}
